\SourceLeaf{396}%
\Needspace{8\baselineskip}
\begin{center}\bfseries [3.]\end{center}
\noindent\hspace*{1.2em}Für das Mass der Seitenkrümmung der Linie, für die $\theta$ constant ist, [stellen wir folgende Betrachtung an.

In diesem Falle sei $t=u$, also
\[
t'=1,\qquad \theta'=0.
\]
Alsdann wird der Reihe nach:]
\begin{equation}\tag*{[$5'$]}
aa+bb+cc=ss,\qquad ss'=aa'+bb'+cc'
\end{equation}
\begin{equation}\tag*{[$4'$]}
\left\{\begin{aligned}s\xi&=a\\s\eta&=b\\s\zeta&=c\end{aligned}\right.
\end{equation}
\begin{equation}\tag*{[$12'$]}
\left\{\begin{aligned}
\tau s\xi'+\xi s'&=a'\\
\tau s\eta'+\eta s'&=b'\\
\tau s\zeta'+\zeta s'&=c'
\end{aligned}\right.
\end{equation}
\begin{equation}\tag*{[$8'$]}
\left\{\begin{aligned}
\eta c'-\zeta b'&=\tau s\xi''\\
\zeta a'-\xi c'&=\tau s\eta''\\
\xi b'-\eta a'&=\tau s\zeta''.
\end{aligned}\right.
\end{equation}
[Werden die letzten drei Gleichungen der Reihe nach mit $\sigma sX,\ \sigma sY,\ \sigma sZ$ multiplicirt und addirt, so ergibt sich wegen 9:
\[
\begin{aligned}
\sigma\tau ss\sin\psi=\sigma s\{&a'(Y\zeta-Z\eta)+b'(Z\xi-X\zeta)\\
&+c'(X\eta-Y\xi)\}.
\end{aligned}
\]
Es ist aber:
\[
\sigma s(Y\zeta-Z\eta)=(c\alpha-a\gamma)c-(a\beta-b\alpha)b,
\]
also]
\[
\sigma s(Y\zeta-Z\eta)=(aa+bb+cc)\alpha-a(a\alpha+b\beta+c\gamma)
\]
[oder
\[
=P\alpha-Qa.
\]
Mithin wird:]
\[
\sigma s^3.\frac{\tau}{s}\sin\psi=P(\alpha a'+\beta b'+\gamma c')-Q(aa'+bb'+cc').
\]
